\ifx\litemode\undefined\documentclass{article}\fi
\usepackage[margin=3cm]{geometry}

\usepackage{amsmath,amssymb}
\usepackage{mathtools}
\usepackage{algorithm}
\usepackage[noend]{algpseudocode}
\usepackage{amsthm}
\usepackage{boxedminipage}
\usepackage{float}
\usepackage{hyperref}
\usepackage{cleveref}

% Algorithmic
\algrenewcommand\textproc{}

% Custom commands
\newcommand{\target}{\mathsf{target}}
\newcommand{\jtargets}{\mathsf{jtargets}}
\newcommand{\ntargets}{\mathsf{ntargets}}
\newcommand{\kind}{\mathsf{kind}}
\newcommand{\height}{\mathsf{height}}
\newcommand{\finalize}{\mathsf{finalize}}
\newcommand{\validator}{\mathsf{validator}}
\newcommand{\justify}{\mathsf{justify}}
\newcommand{\notarize}{\mathsf{notarize}}
\newcommand{\cnt}{\mathsf{count}}
\newcommand{\parent}{\mathsf{parent}}
\newcommand{\sentfinalize}{\mathsf{sent\_finalize}}
\newcommand{\sentjustify}{\mathsf{sent\_justify}}
\newcommand{\sentnotarize}{\mathsf{sent\_notarize}}
\newcommand{\getKDeep}{\mathsf{getKDeep}}
\newcommand{\getConfirmed}{\mathsf{getConfirmed}}
\newcommand{\true}{\mathsf{true}}

% Theorem environments
\theoremstyle{plain}
\newtheorem{theorem}{Theorem}
\newtheorem{lemma}{Lemma}

\theoremstyle{definition}
\newtheorem{definition}[theorem]{Definition}

\theoremstyle{remark}
\newtheorem{remark}[theorem]{Remark}

\title{Finalization as Notarization}

\IfFileExists{lite_full_setup.tex}{% Shared lite/full setup for protocol-spec .tex files.
%
% Files under full/ that include this snippet compile in two modes:
%   - Full (default): all content rendered.
%   - Lite: when \litemode is defined before \documentclass (typically by a
%     wrapper .tex file at the repo root), lemmas, corollaries, propositions,
%     properties, proofs, remarks, assumptions, and conditions are dropped
%     via the comment package. Algorithms, theorem statements, definitions,
%     and surrounding prose remain.
%
% Place \input of this file in each spec's preamble AFTER all \newtheorem
% declarations (or just before \begin{document}). The exclusions use
% \@ifundefined so it's safe even when a file defines only a subset.
%
% Path resolution: \IfFileExists tries the local copy first (when cwd is
% full/), then falls back to full/lite_full_setup.tex (when cwd is the repo
% root, e.g. while compiling a top-level _lite.tex wrapper).
\newif\iffull
\ifx\litemode\undefined\fulltrue\else\fullfalse\fi
\iffull\else
  \usepackage{comment}
  \makeatletter
  \@ifundefined{lemma}{}{\excludecomment{lemma}}
  \@ifundefined{corollary}{}{\excludecomment{corollary}}
  \@ifundefined{proposition}{}{\excludecomment{proposition}}
  \@ifundefined{property}{}{\excludecomment{property}}
  \@ifundefined{proof}{}{\excludecomment{proof}}
  \@ifundefined{remark}{}{\excludecomment{remark}}
  \@ifundefined{assumption}{}{\excludecomment{assumption}}
  \@ifundefined{condition}{}{\excludecomment{condition}}
  \makeatother
\fi
}{% Shared lite/full setup for protocol-spec .tex files.
%
% Files under full/ that include this snippet compile in two modes:
%   - Full (default): all content rendered.
%   - Lite: when \litemode is defined before \documentclass (typically by a
%     wrapper .tex file at the repo root), lemmas, corollaries, propositions,
%     properties, proofs, remarks, assumptions, and conditions are dropped
%     via the comment package. Algorithms, theorem statements, definitions,
%     and surrounding prose remain.
%
% Place \input of this file in each spec's preamble AFTER all \newtheorem
% declarations (or just before \begin{document}). The exclusions use
% \@ifundefined so it's safe even when a file defines only a subset.
%
% Path resolution: \IfFileExists tries the local copy first (when cwd is
% full/), then falls back to full/lite_full_setup.tex (when cwd is the repo
% root, e.g. while compiling a top-level _lite.tex wrapper).
\newif\iffull
\ifx\litemode\undefined\fulltrue\else\fullfalse\fi
\iffull\else
  \usepackage{comment}
  \makeatletter
  \@ifundefined{lemma}{}{\excludecomment{lemma}}
  \@ifundefined{corollary}{}{\excludecomment{corollary}}
  \@ifundefined{proposition}{}{\excludecomment{proposition}}
  \@ifundefined{property}{}{\excludecomment{property}}
  \@ifundefined{proof}{}{\excludecomment{proof}}
  \@ifundefined{remark}{}{\excludecomment{remark}}
  \@ifundefined{assumption}{}{\excludecomment{assumption}}
  \@ifundefined{condition}{}{\excludecomment{condition}}
  \makeatother
\fi
}

\begin{document}
\maketitle

\section{Model and definitions}

We consider a system of~$n$ validators of which at most~$f$ are adversarial, where $n \geq 3f + 1$.
A \emph{quorum} is any set of at least $2n/3$ validators.
Any two quorums overlap in at least $n/3 + 1 > f$ validators, so their intersection contains at least one honest validator.

\begin{definition}[Block]\label{def:block}
A \emph{block} has a \emph{parent} (another block) and contains a set of votes (Definition~\ref{def:vote}).
The \emph{genesis block}~$B_{\mathrm{genesis}}$ has no parent.
\end{definition}

\begin{definition}[Chain and divergence]\label{def:chain}
A \emph{chain} is a sequence of blocks forming a path from the genesis block, where each block's parent is its predecessor in the sequence.
Since a chain is uniquely determined by its last block, we identify chains with their last block and use the two interchangeably.
Two chains \emph{conflict} if neither is a prefix of the other.
\end{definition}

We write $B \preceq B'$ when $B$ is an ancestor of~$B'$ (or $B = B'$), and $B \sim B'$ when $B$ and $B'$ are \emph{compatible}: $B \preceq B'$ or $B' \preceq B$.
For compatible blocks, we write $B > B'$ to mean $B \succ B'$ (i.e., $B$ is a strict descendant of~$B'$); we say $B$ is \emph{longer} or \emph{greater} than~$B'$.

\begin{definition}[Vote]\label{def:vote}
A \emph{vote} is a signed tuple $(\validator,\; \height,\; \target,\; \kind)$ where $\validator$ identifies the signer, $\height$ is a height, $\target$ is a block, and $\kind \in \{\notarize, \justify,\, \finalize\}$ distinguishes the vote type.
A vote included on a chain is verified to reference a target that actually exists on that chain.
\end{definition}

\begin{definition}[Slashing condition]\label{def:slashing}
Given two votes~$a$, $b$ from the same validator, the validator is \emph{slashable} if either of the following conditions hold
\begin{enumerate}
  \item \(
    a.\kind \in \{\justify,\finalize\}
    \;\;\wedge\;\;
    a.\kind = b.\kind
    \;\;\wedge\;\;
    a.\height = b.\height
    \;\;\wedge\;\;
    a.\target \neq b.\target.
  \)

  \item \(
    a.\kind = \finalize
    \;\;\wedge\;\;
    b.\kind = \notarize
    \;\;\wedge\;\;
    a.\height = b.\height.
  \)
\end{enumerate}

That is: (1)~a validator cannot cast two justify (resp.\ finalize) votes at the same height with different targets, and (2)~a validator cannot cast both a finalize vote and a notarize vote at the same height.
We refer to a slashable validator also as an \emph{equivocator}.

These are the \emph{only} slashing conditions.
There is no double-target condition for notarize votes: validators may cast multiple notarize votes at the same height without slashing risk.
\end{definition}

\section{Algorithms}

\emph{Convention.} Maps accessed at an uninitialized key return $\bot$ by default.
Every block satisfies $B \succeq \bot$.

\begin{figure}[H]
\caption{State machine}\label{alg:state-machine}
\begin{center}
\begin{boxedminipage}{\textwidth}
\begin{algorithmic}[1]

\Function{$\textsf{processBlock}$}{$s,\; B$}
  \For{each vote $v$ in $B$} $s \gets \textsf{processVote}(s,\, v,\, B)$ \EndFor
  \State \Return $\textsf{processHeight}(s,\, B)$
\EndFunction
\vspace{3pt}

\Function{$\textsf{processVote}$}{$(h,\, \jtargets,\, \ntargets,\, T^\mathsf{min},\, J,\, h_j,\, F,\, P,\, T^\mathsf{max}),\; v,\; B$}
  \State $i \gets v.\validator$
  \If{$v.\height = h$ and $v.\target \preceq B$}
    \If{$v.\kind = \justify$}
      \If{$\jtargets[h][i] = \bot$}
        \State $\jtargets[h][i] \gets v.\target$
      \EndIf
    \Else  \Comment{Both notarize and finalize votes}
      \If{$\ntargets[i] = \bot$ or $v.\target \preceq \ntargets[i]$}
        \State $\ntargets[i] \gets v.\target$
      \EndIf
      \If{$T^\mathsf{min} = \bot$ or $v.\target \preceq T^\mathsf{min}$}
        \State $T^\mathsf{min} \gets v.\target$
      \EndIf
    \EndIf
  \EndIf
  \If{$v.\height = h_j$ and $v.\kind = \finalize$ and $\jtargets[h_j][i] \neq \bot$ and $v.\target \preceq \jtargets[h_j][i]$ and $J \preceq v.\target \preceq B$} \Comment{Finalize votes only}
    \State $P \gets P \cup \{i\}$
  \EndIf
  \State \Return $(h,\, \jtargets,\, \ntargets,\, T^\mathsf{min},\, J,\, h_j,\, F,\, P,\, T^\mathsf{max})$
\EndFunction
\vspace{3pt}

\Function{$\textsf{processHeight}$}{$(h,\, \jtargets,\, \ntargets,\, T^\mathsf{min},\, J,\, h_j,\, F,\, P,\, T^\mathsf{max}),\; B$}
  \If{$|P| \geq 2n/3$} \Comment{Finality}
    \State $F \gets J$
  \EndIf
  \State $\cnt \gets 0$ \Comment{Justification: first pass, walk from $B$ down to $T^\mathsf{max}$}
  \For{$B'$ from $B$ down to $T^\mathsf{max}$}
    \State $\cnt \gets \cnt + |\{i : \jtargets[h][i] = B'\}|$
    \If{$\cnt \geq 2n/3$} \Comment{Height advances; $B' \succeq T^\mathsf{max} \succeq J$ by walk range}
      \State $\cnt \gets 0$ \Comment{Second pass: find tightest $T^\mathsf{max\_temp}$}
      \For{$B''$ from $B'$ up to $B$}
        \State $\cnt \gets \cnt + |\{i : \jtargets[h][i] = B''\}|$
        \If{$\cnt \geq 2n/3$}
          \State $T^\mathsf{max\_temp}\gets B''$
          \State \textbf{break}
        \EndIf
      \EndFor
      \State \textbf{delete} $\jtargets[h_j]$
      \State $J \gets B'$, \; $h_j \gets h$, \; $P \gets \emptyset$
      \If{$T^\mathsf{max\_temp} \succeq T^\mathsf{max}$}
        \State $T^\mathsf{max}\gets T^\mathsf{max\_temp}$
      \EndIf
      \State $h \gets h + 1$
      \State \textbf{delete} $\ntargets$, $T^\mathsf{min}$
      \State \Return $(h,\, \jtargets,\, \ntargets,\, T^\mathsf{min},\, J,\, h_j,\, F,\, P,\, T^\mathsf{max})$
    \EndIf
  \EndFor
  \State $\cnt \gets 0$ \Comment{Notarization: walk from $T^\mathsf{min}$ toward $B$}
  \For{$B'$ from $T^\mathsf{min}$ up to $B$}
    \State $\cnt \gets \cnt + |\{i : \ntargets[i] = B'\}|$
    \If{$\cnt \geq 2n/3$} \Comment{Height advances}
      \If{$B' \succeq T^\mathsf{max}$}
        \State $T^\mathsf{max}\gets B'$
      \Else
        \State \textbf{delete} $\jtargets[h]$
      \EndIf
      \State $h \gets h + 1$
      \State \textbf{delete} $\ntargets$, $T^\mathsf{min}$
      \State \Return $(h,\, \jtargets,\, \ntargets,\, T^\mathsf{min},\, J,\, h_j,\, F,\, P,\, T^\mathsf{max})$
    \EndIf
  \EndFor
  \State \Return $(h,\, \jtargets,\, \ntargets,\, T^\mathsf{min},\, J,\, h_j,\, F,\, P,\, T^\mathsf{max})$
\EndFunction

\end{algorithmic}
\end{boxedminipage}
\end{center}
\end{figure}

\section{Accountable safety}

\begin{lemma}[Height progression]\label{lem:heightprog}
For any chain~$Y$ and $h' \in [1,\sigma[Y].h-1]$, there exists a block $B \preceq Y$ such that $\sigma[B].h = h'$.
\end{lemma}

\begin{proof}
The $\textsf{processBlock}$ function computes the post-state of each block from its parent's post-state via $\textsf{processHeight}$, which increments~$h$ by at most~$1$ (in the justification or notarize branches) or leaves it unchanged.
Therefore, as blocks are processed along the chain from genesis to~$Y$, the height~$h$ increases one step at a time.
Since $\sigma[B_\mathrm{genesis}].h = 0$ and $\sigma[Y].h > h'$, there must exist a block $B \preceq Y$ at which the post-state height first reaches~$h'$, i.e., $\sigma[B].h = h'$.
\end{proof}

\begin{lemma}[Justification compatibility per height]\label{lem:uniqueness}
Unless $\geq n/3$ validators are slashable: only blocks on the same chain can be justified at any given height.
More precisely, if block~$T_1$ is justified at height~$h$ on chain~$X$ and block~$T_2$ is justified at height~$h$ on chain~$Y$, then $T_1 \sim T_2$.
\end{lemma}

\begin{proof}
Justification of~$T_1$ at height~$h$ on chain~$X$ requires a quorum~$Q_1$ ($|Q_1| \geq 2n/3$) where each member~$i$ has $\jtargets[h][i] \succeq T_1$ on~$X$.
Likewise, justification of~$T_2$ at height~$h$ on chain~$Y$ requires a quorum~$Q_2$ ($|Q_2| \geq 2n/3$) where each member has $\jtargets[h][i] \succeq T_2$ on~$Y$.

Suppose for contradiction that $T_1 \not\sim T_2$, i.e., they lie on different branches past some fork point~$G$.
Every target in~$Q_1$ is a descendant of~$T_1$, hence strictly on branch~$X$ past~$G$; similarly every target in~$Q_2$ is strictly on branch~$Y$ past~$G$.

By quorum intersection, $|Q_1 \cap Q_2| \geq n/3$.
Each validator~$i$ in $Q_1 \cap Q_2$ has $\jtargets[h][i]$ on branch~$X$ (for~$Q_1$) and on branch~$Y$ (for~$Q_2$).
Since only justify votes populate $\jtargets$ and $\textsf{processVote}$ requires $v.\target \preceq B$ (the target must lie on the chain being processed), a single justify vote cannot contribute to both chains past the fork.
Therefore~$i$ must have signed two justify votes at height~$h$ with different targets, which is slashable by condition~1.
Since all $\geq n/3$ validators in the intersection are slashable, this contradicts the assumption.
\end{proof}

\begin{lemma}[Notarization respects finalization]\label{lem:finexclude}
Unless $\geq n/3$ validators are slashable: if $C$ is finalized at height~$h$ and block~$B$ is notarized at height~$h$ on some chain, then $C \preceq B$.
\end{lemma}

\begin{proof}
Finalization of~$C$ at height~$h$ requires a quorum~$Q_F$ ($|Q_F| \geq 2n/3$) where each member~$i$ signed a finalize vote with $\height = h$ and $\target \succeq C$.
By slashing condition~2, no member of~$Q_F$ can have signed a notarize vote at height~$h$ without being slashable.
By slashing condition~1, each member has a unique finalize target $D_i \succeq C$, so $\ntargets[i] = D_i$.

Notarization of~$B$ at height~$h$ means the cumulative count in the notarize walk reached $2n/3$ at~$B$.
The walk proceeds from $T^\mathsf{min}$ toward~$B$, and $\ntargets[i]$ records each validator's shallowest notarize/finalize target.
Let $Q_N$ be the set of validators counted at or before~$B$, i.e., $Q_N = \{i : \ntargets[i] \preceq B\}$, with $|Q_N| \geq 2n/3$.
By quorum intersection, $|Q_F \cap Q_N| \geq n/3$.

Consider a non-slashable validator~$i$ in $Q_F \cap Q_N$.
From~$Q_F$: $\ntargets[i] = D_i \succeq C$.
From~$Q_N$: $\ntargets[i] = D_i \preceq B$.
Combined: $C \preceq D_i \preceq B$, hence $C \preceq B$.
\end{proof}


\begin{lemma}[Justification respects finalization]\label{lem:tbd}
Unless $\geq n/3$ validators are slashable: if $C$ is finalized at height~$h$ and block~$B$ is justified at height~$h$ on some chain $X$, then $C \preceq X$.
\end{lemma}

\begin{proof}
Finalization of~$C$ at height~$h$ requires a quorum~$Q_F$ ($|Q_F| \geq 2n/3$) where each member~$i$ signed a finalize vote with target $D_i \succeq C$.
By the finalization condition in $\textsf{processVote}$, each member also has a justify vote at height~$h$ with target $\jtargets[h][i] \succeq D_i$.

Justification of~$B$ at height~$h$ on chain~$X$ requires a quorum~$Q_J$ ($|Q_J| \geq 2n/3$) where each member has $\jtargets[h][i] \succeq B$ on~$X$, so in particular $\jtargets[h][i] \preceq X$.

By quorum intersection, $|Q_F \cap Q_J| \geq n/3$.
A non-slashable validator~$i$ in $Q_F \cap Q_J$ has a single justify target at height~$h$ (condition~1), serving both quorums.
From~$Q_F$: $C \preceq D_i \preceq \jtargets[h][i]$.
From~$Q_J$: $\jtargets[h][i] \preceq X$.
Hence $C \preceq X$.
\end{proof}

\begin{lemma}[Any chain past a finalized height contains the finalized block]\label{lem:mainsafety}
Unless $\geq n/3$ validators are slashable: if $C$ is finalized at height~$h$, then any chain~$Y$ that progresses past height~$h$ satisfies $C \preceq Y$.
\end{lemma}

\begin{proof}
Chain~$Y$ advances through height~$h$ via either justification or notarization.
If via justification: Lemma~\ref{lem:tbd} gives $C \preceq Y$.
If via notarization of some block~$B$: Lemma~\ref{lem:finexclude} gives $C \preceq B$, and $B \preceq Y$ (since $B$ is on chain~$Y$), so $C \preceq Y$.
\end{proof}

\begin{lemma}[Finalized blocks form a chain]\label{lem:finchain}
Unless $\geq n/3$ validators are slashable: if $C$ is finalized at height~$h$ and $C'$ is finalized at height~$h' \geq h$, then $C \sim C'$.
\end{lemma}

\begin{proof}
If $h = h'$: both $C$ and $C'$ are justified at height~$h$ (finalization requires prior justification).
By Lemma~\ref{lem:uniqueness}, $C \sim C'$.

If $h' > h$: let $X'$ be the chain on which $C'$ is finalized.
Chain~$X'$ has progressed past height~$h$ (since $h' > h$).
By Lemma~\ref{lem:mainsafety}, $C \preceq X'$.
Since $C' \preceq X'$ (it is on chain~$X'$), both $C$ and $C'$ lie on the same chain, hence $C \sim C'$.
\end{proof}

\begin{theorem}[Accountable safety]\label{thm:safety}
Unless $\geq n/3$ validators are slashable, no two conflicting blocks can be finalized.
\end{theorem}

\begin{proof}
Suppose $C$ is finalized at height~$h$ and $C'$ at height~$h'$, with $h \leq h'$ without loss of generality.
By Lemma~\ref{lem:finchain}, $C \sim C'$, so $C$ and $C'$ are not conflicting.
\end{proof}

\begin{figure}[H]
\caption{Store and voting}\label{alg:store}
\begin{center}
\begin{boxedminipage}{\textwidth}
\begin{algorithmic}[1]

\State \textbf{Store:} $J_s \gets B_{\mathrm{genesis}}$, \; $F_s \gets B_{\mathrm{genesis}}$, \; $\mathcal{T} \gets \{B_{\mathrm{genesis}}\}$, \; $C \gets \emptyset$
\State $\sigma[B_{\mathrm{genesis}}] \gets (0,\; \bot^n,\; \bot^n,\; \bot,\; B_{\mathrm{genesis}},\; 0,\; B_{\mathrm{genesis}},\; \emptyset,\; \bot)$
\State \textbf{Per-validator:} $\sentjustify \gets \bot^*$, \; $\sentnotarize \gets \bot^*$, \; $\sentfinalize \gets \bot^*$
\vspace{3pt}

\Function{$\textsf{onBlock}$}{$B$}
  \State \textbf{assert} $B.\parent \in \mathcal{T}$ \Comment{Parent already accepted}
  \State \textbf{assert} $F_s \preceq B$ \Comment{$B$ descends from finalized}
  \State $\mathcal{T} \gets \mathcal{T} \cup \{B\}$
  \State $\sigma[B] \gets \textsf{processBlock}(\sigma[B.\parent],\; B)$
  \State $C \gets C \cup \{(\sigma[B].J,\; \sigma[B].h_j)\}$
  \State $J_s \gets \textsf{computeJustified}()$
  \State $\textsf{updateFinalized}(\sigma[B].F)$
\EndFunction
\vspace{3pt}

\Function{$\textsf{computeJustified}$}{}
  % \State $J^* \gets F_s$
  % \While{$\exists\, (J',\, h') \in C$ with $J' \succ J^*$} \Comment{Walk toward descendants}
  \State \Return $\arg\max_{(J',\, h') \in C,\; J' \succeq F_s} (h',\, J')$
  % \EndWhile
  % \State \Return $J^*$
\EndFunction
\vspace{3pt}

\Function{$\textsf{updateFinalized}$}{$F_B$}
  \If{$F_B \succ F_s$ and $F_B \preceq J_s$}
    \State $F_s \gets F_B$
  \EndIf
\EndFunction

\Function{$\textsf{onVote}$}{}
  \State $h^\mathsf{can}\gets \sigma[\textsf{getHead}()].h$
  \If{$\lnot\, \sentfinalize[h^\mathsf{can}]$}
    \If{$\lnot\, \sentjustify[h^\mathsf{can}]$}
      \State \textbf{send} $(i, h^\mathsf{can}, \getConfirmed(), \justify)$
      \State $\sentjustify[h^\mathsf{can}]\gets \true$
    \Else
      \State \textbf{send} $(i, h^\mathsf{can}, \getKDeep(), \notarize)$
      \State $\sentnotarize[h^\mathsf{can}]\gets \true$
    \EndIf
  \EndIf
  \State $(J_j, h_j) \gets \arg\max_{(J',h')\in C}(h', J')$
  \If{$\lnot\, \sentfinalize[h_j] \land \lnot\, \sentnotarize[h_j]$}
      \State \textbf{send} $(i, h_j, J_j, \finalize)$
      \State $\sentfinalize[h_j]\gets \true$
  \EndIf
\EndFunction

\end{algorithmic}
\end{boxedminipage}
\end{center}
\end{figure}

\section{Inactivity Leak}

\begin{lemma}[Justifications are compatible with finalized blocks]\label{lem:justcompat}
Unless $\geq n/3$ validators are slashable: if $C$ is finalized at height~$h$, then any block justified at height $\geq h$ on any chain is compatible with~$C$.
\end{lemma}

\begin{proof}
Let $B$ be justified at height~$h'$ on chain~$Y$, with $h' \geq h$.
If $h' = h$: both $C$ and $B$ are justified at height~$h$, so $C \sim B$ by Lemma~\ref{lem:uniqueness}.
If $h' > h$: chain~$Y$ has progressed past height~$h$ (since it reached height~$h'$).
By Lemma~\ref{lem:mainsafety}, $C \preceq Y$.
Since $B \preceq Y$ (it is on chain~$Y$), both $C$ and $B$ lie on the same chain, hence $C \sim B$.
\end{proof}

\begin{lemma}[$J$ monotonicity and the $T^\mathsf{max}$ invariant]\label{lem:jmono}
On any single chain, $J$ is monotonically non-decreasing: each justification sets $J$ to a descendant of (or equal to) the previous~$J$.
Moreover, the invariant $T^\mathsf{max} \succeq J$ holds at all times.
\end{lemma}

\begin{proof}
The justification walk iterates from~$B$ down to~$T^\mathsf{max}$.
When the quorum is reached at block~$B'$, we have $B' \succeq T^\mathsf{max}$ (since $B'$ is in the walk range).
The algorithm then sets $J \gets B'$ and updates $T^\mathsf{max} \gets T^\mathsf{max\_temp}$ (when $T^\mathsf{max\_temp} \succeq T^\mathsf{max}$, which holds since $T^\mathsf{max\_temp}$ is also in the walk range).
Since $T^\mathsf{max\_temp} \succeq B' = J$, the invariant $T^\mathsf{max} \succeq J$ is established.

Notarization may only increase $T^\mathsf{max}$ and does not modify~$J$, so the invariant is preserved.

For monotonicity: at the next justification, the walk range is $[T^\mathsf{max}, B]$, so $B' \succeq T^\mathsf{max} \succeq J_{\mathrm{old}}$.
Therefore $J_{\mathrm{new}} = B' \succeq J_{\mathrm{old}}$.
\end{proof}

\begin{lemma}\label{lem:witnessed}
  For any chain $B$ and $h' \in [1, \sigma[B].h_j-1]$, there exists a set of votes, each with target $\preceq \sigma[B].J$, that either notarizes or justifies a block $\preceq B$ at height $h'$.
\end{lemma}

\begin{proof}
Let $J = \sigma[B].J$ and $h_j = \sigma[B].h_j$.
Since $h_j \leq \sigma[B].h - 1$, we have $h' \leq h_j - 1 \leq \sigma[B].h - 2$, so $h' \in [1, \sigma[B].h - 1]$.
By Lemma~\ref{lem:heightprog}, there exists a block $B' \preceq B$ with $\sigma[B'].h = h'$.
Height advanced from~$h'$ to~$h'+1$ when $\textsf{processHeight}$ was called on~$B'$, via either the justification or notarization walk.

\emph{Justification case.}
The justify walk runs from~$B'$ down to~$T^\mathsf{max}$, so every counted vote has target in $[T^\mathsf{max},\, B']$.
The deepest target in this range is $T^\mathsf{max\_temp}$, and $T^\mathsf{max}$ is updated to $T^\mathsf{max\_temp}$ when the quorum is reached.
By Lemma~\ref{lem:jmono}, $T^\mathsf{max}$ only increases and $J \succeq T^\mathsf{max}$ at all times, so $J \succeq T^\mathsf{max\_temp}$.
Therefore every vote target at height~$h'$ satisfies $\mathsf{target} \preceq T^\mathsf{max\_temp} \preceq J$.
The justified block is $\preceq B' \preceq B$.

\emph{Notarization case.}
The notarize walk runs from~$T^\mathsf{min}$ up to~$B'$, and the quorum is reached at some block~$B''$.
Each counted validator has $\ntargets[i] \preceq B''$.
After the event, $T^\mathsf{max} \gets \max(T^\mathsf{max},\, B'')$.
By Lemma~\ref{lem:jmono}, $J \succeq T^\mathsf{max}$ at all times: if $B'' \succeq T^\mathsf{max}$, then $T^\mathsf{max}$ is updated to $B''$ and $J \succeq T^\mathsf{max} \succeq B''$; if $B'' \prec T^\mathsf{max}$, then $B'' \prec T^\mathsf{max} \preceq J$.
In either case $B'' \preceq J$, so every vote target satisfies $\mathsf{target} \preceq B'' \preceq J$.
The notarized block $B''$ is $\preceq B' \preceq B$.
\end{proof}

\begin{lemma}[Justification is bounded by the longest target]\label{lem:justbound}
Unless $\geq n/3$ validators are slashable: if a quorum of justify votes at height~$h$ justifies block~$C$, and $T$ is the longest target among those votes, then every block~$C'$ justified at height~$h$ (by any quorum) satisfies $C' \preceq T$.
\end{lemma}

\begin{proof}
Let $Q_1$ ($|Q_1| \geq 2n/3$) be the set of justify votes at height~$h$ that justifies~$C$: each member~$i$ has target $t_i \succeq C$, and $T$ is the longest target, so $t_i \preceq T$ for all $i \in Q_1$.
Let $Q_2$ ($|Q_2| \geq 2n/3$) be the set that justifies~$C'$: each member has target $\succeq C'$.

By quorum intersection, $|Q_1 \cap Q_2| \geq n/3$.
A non-slashable validator~$i$ in $Q_1 \cap Q_2$ has a unique justify target at height~$h$ (condition~1).
From~$Q_1$: $t_i \preceq T$.
From~$Q_2$: $t_i \succeq C'$.
Hence $C' \preceq t_i \preceq T$.
\end{proof}




\begin{lemma}\label{lem:upgrade-helper}
  Assume $< n/3$ of validators are slashable.
  Let $(J_j, h_j) \in C$.
  If $F$ is finalized at height~$h_F < h_j$ and $(F, h_F) \in C$, then $J_j \succeq F$.
\end{lemma}

\begin{proof}
Since $F$ is finalized at $h_F$ and $J_j$ is justified at $h_j > h_F$, Lemma~\ref{lem:justcompat} gives $J_j \sim F$.
If $J_j \succeq F$ we are done, so assume for contradiction that $J_j \prec F$.

Let $B$ be the block on whose chain $(J_j, h_j)$ was produced, i.e., $\sigma[B].J = J_j$ and $\sigma[B].h_j = h_j$.
Since $h_F < h_j$, Lemma~\ref{lem:witnessed} yields a set of votes at height~$h_F$, each with target $\preceq J_j$, that either justifies or notarizes a block $\preceq B$.

\emph{Justification case.}
The witness is a quorum of justify votes at $h_F$ with longest target $T \preceq J_j$.
By Lemma~\ref{lem:justbound}, every block justified at $h_F$ satisfies $C' \preceq T$.
Since finalization of~$F$ requires prior justification at~$h_F$: $F \preceq T \preceq J_j$, contradicting $J_j \prec F$.

\emph{Notarization case.}
The witness notarizes a block~$B'$ at height~$h_F$ with $B' \preceq J_j$.
By Lemma~\ref{lem:finexclude}: $F \preceq B' \preceq J_j$, again contradicting $J_j \prec F$.
\end{proof}

\begin{lemma}[Upgrade property]\label{lem:upgrade}
  Assume $< n/3$ of validators are slashable.
  Let $(J_j, h_j) = \arg\max_{(J',h')\in C \land J' \succeq F_s}(h', J')$.
  If $F$ is finalized at height~$h_F$ and $(F, h_F) \in C$, then $J_j \succeq F$.
\end{lemma}

\begin{proof}
By accountable safety (Theorem~\ref{thm:safety}), $F \sim F_s$.

\emph{Case $F \preceq F_s$.}
$J_j \succeq F_s \succeq F$.

\emph{Case $F \succ F_s$.}
Then $(F, h_F) \in C$ with $F \succeq F_s$, so it is a candidate for the $\arg\max$.
If $h_F < h_j$: Lemma~\ref{lem:upgrade-helper} gives $J_j \succeq F$.
If $h_F \geq h_j$: since $(J_j, h_j)$ is the $\arg\max$ and $(F, h_F)$ is a candidate with $h_F \geq h_j$, we must have $h_F = h_j$ and $J_j \geq F$ in the tiebreak ordering.
By Lemma~\ref{lem:uniqueness}, $F \sim J_j$ (both justified at height~$h_j$), so $J_j \succeq F$.
\end{proof}

\begin{lemma}[Honest finalize targets respect the store]\label{lem:honest-fin}
  Assume $< n/3$ of validators are slashable.
  Let $(J_j, h_j) = \arg\max_{(J',h')\in C \land J' \succeq F_s}(h', J')$ in the view of a validator~$i$.
  If $i$ is honest and casts a $\finalize$ vote at height~$h_j$ with target~$F$, then $F \preceq J_j$.
\end{lemma}

\begin{proof}
From $\textsf{onVote}$, the honest validator computes $(F, h_j) = \arg\max_{(J',h')\in C}(h', J')$ and sends a finalize vote with target~$F$ at height~$h_j$.
In particular, $(F, h_j) \in C$, so $F$ is justified at height~$h_j$.
Since $J_j$ is also justified at height~$h_j$, Lemma~\ref{lem:uniqueness} gives $F \sim J_j$.

If $F \succeq F_s$: then $(F, h_j)$ is a candidate for the restricted $\arg\max$.
Since $(J_j, h_j)$ is the winner at the same height, $J_j \geq F$ in the tiebreak ordering, and $F \sim J_j$ gives $F \preceq J_j$.

If $F \not\succeq F_s$: suppose for contradiction that $F \succ J_j$.
Then $F \succ J_j \succeq F_s$, so $F \succeq F_s$---a contradiction.
Therefore $F \preceq J_j$.
\end{proof}

% \subsection{WIP}

\begin{lemma}
  % Assume $< n/3$ of validators are slashable.
  Assume an honest validator $i$.
  Let $h_F$ be such that $(F_s, h_F) \in C$ for $i$.
  Let $(J_j, h_j) = \arg\max_{(J',h')\in C \land J' \succeq F_s}(h', J')$ in the view of a validator $i$.
  Then $h_j \geq h_F$.
\end{lemma}

\begin{proof}
$(F_s, h_F) \in C$ with $F_s \succeq F_s$, so it is a candidate for the $\arg\max$.
Since $(J_j, h_j)$ is the maximizer, $(h_j, J_j) \geq (h_F, F_s)$ lexicographically, hence $h_j \geq h_F$.
\end{proof}

% \begin{lemma}
%   Assume $< n/3$ of validators are slashable.
%   Assume an honest validator $i$.
%   Let $h_F$ be such that $(F_s, h_F) \in C$ for $i$.
%   Let $(J_j, h_j) = \arg\max_{(J',h')\in C \land J' \succeq F_s}(h', J')$ in the view of a validator $i$.
%   Take a chain $B$ extending $J_j$.
%   Take $h' \in [1, \max(h_F,h_j)-1]$.
%   there exists a set of votes, each with target $\preceq B$, that either notarizes or justifies a block $\preceq B$ at height $h'$.
% \end{lemma}

% \begin{proof}
% Since $B \succeq J_j \succeq F_s$, both $J_j$ and $F_s$ lie on chain~$B$.
% Let $(J^*, h^*)$ be whichever of $(J_j, h_j)$ and $(F_s, h_F)$ has the larger height, so $h^* = \max(h_F, h_j)$.
% Since $(J^*, h^*) \in C$ and $J^* \preceq J_j \preceq B$, Lemma~\ref{lem:witnessed} applied to the chain that produced $(J^*, h^*)$ yields, for each $h' \in [1, h^* - 1]$, a witness with targets $\preceq J^* \preceq B$ that justifies or notarizes a block $\preceq B$.
% \end{proof}


% \begin{lemma}
%   Assume $< n/3$ of validators are slashable.
%   Assume an honest validator $i$.
%   Let $h_F$ be such that $(F_s, h_F) \in C$ for $i$.
%   Let $(J_j, h_j) = \arg\max_{(J',h')\in C \land J' \succeq F_s}(h', J')$ in the view of a validator $i$.
%   Take a chain $B$ extending $J_j$.
%   If $i$ casts a $\finalize$ vote at height~$\max(h_j,h_F)$ with target~$F$, then $F \preceq B$.
% \end{lemma}

% \begin{proof}
%   Say $h_j \geq h_F$, then use \Cref{lem:honest-fin}.

%   If $h_j < h_F$, then 
% \end{proof}

\begin{lemma}
  Assume $< n/3$ of validators are slashable.
  Assume an honest validator $i$.
  Let $(J_j, h_j) = \arg\max_{(J',h')\in C \land J' \succeq F_s}(h', J')$ in the view of a validator $i$.
  Validator $i$ has never sent a finalize message for height $>h_j$.
\end{lemma}

\begin{proof}
Suppose for contradiction that $i$ sent a finalize vote at height~$h_B > h_j$ with target~$B_F$.
From $\textsf{onVote}$, $(B_F, h_B) = \arg\max_{(J',h')\in C}(h', J')$, so $(B_F, h_B) \in C$.

Let $h_F$ be such that $(F_s, h_F) \in C$.
By a previous lemma, $h_j \geq h_F$, so $h_B > h_j \geq h_F$.
Since $F_s$ is finalized at height~$h_F$ and $h_F < h_B$, Lemma~\ref{lem:upgrade-helper} (applied to the entry $(B_F, h_B) \in C$) gives $B_F \succeq F_s$.

But then $(B_F, h_B) \in C$ with $B_F \succeq F_s$ and $h_B > h_j$, contradicting $(J_j, h_j)$ being $\arg\max_{J' \succeq F_s} C$.
\end{proof}

% \begin{lemma}
%   Assume $< n/3$ of validators are slashable.
%   Pick any $(J^m,h^m) \in \{(J',h')\in C, \forall (J'',h'')\in C, h'' \leq h'\}$.
%   For any $(J',h')\in C$ such that $h' < h^m$ and $J' \sim J^m$, $J' \preceq J^m$.
% \end{lemma}


% \begin{lemma}
%   Assume $< n/3$ of validators are slashable.
%   Let $(J_j, h_j) \gets \arg\max_{(J',h')\in C \land J' \succeq F_s}(h', J')$.
%   If $F$ is finalized at height~$h_F$, $(F, h_F) \in C$, then produces $J_j \succeq F$.
% \end{lemma}

% \begin{proof}
%   Intuition.
%   By  a previous Lemma, $J' \sim F$.
%   $h_j \geq h_F$.
%   If 
% \end{proof}

% \begin{lemma}
  
% \end{lemma}

\section{Finalize via longest common prefix}

\Cref{alg:state-machine-finprefix} modifies \Cref{alg:state-machine} in two ways.

\paragraph{Change 1: finalize via walk.}
Instead of $F \gets J$ when $|P| \geq 2n/3$, the algorithm walks the finalize targets top-down (as justification walks justify targets) and finalizes the longest block~$F^*$ with $2n/3$ support.
$F$ is updated only if $F^* \succ F$ (finalization never regresses).

\paragraph{Change 2: dropping the $J \preceq v.\target$ guard.}
The finalize condition in \Cref{alg:state-machine} requires $J \preceq v.\target$ (the finalize target must be at least as long as~$J$).
This guard is removed in \Cref{alg:state-machine-finprefix}: any finalize target $v.\target \preceq \jtargets[h_j][i]$ on the chain is accepted.

In \Cref{alg:state-machine} the guard is necessary because finalization sets $F \gets J$ unconditionally---without the guard, validators with finalize targets shorter than~$J$ would enter~$P$, reaching $|P| \geq 2n/3$ even when the quorum does not truly support~$J$.
The safety proofs (e.g., Lemma~\ref{lem:finexclude}) assume every finalize voter has target $\succeq F$; removing the guard would violate this.

In \Cref{alg:state-machine-finprefix} the guard is unnecessary: the finalize walk naturally finds the longest block~$F^*$ with genuine $2n/3$ support.
The $\geq 2n/3$ finalize voters counted at~$F^*$ all have $\mathsf{ftargets}[i] \succeq F^*$, so the safety proofs carry over with~$F^*$ in place of~$J$.
The consequence is that $F^*$ may be shorter than~$J$ (if many finalize targets cluster below~$J$), yielding weaker finalization than \Cref{alg:state-machine} in that case.

\label{alg:state-machine-finprefix}
The following changes are applied to \Cref{alg:state-machine}.
The state is extended with $\mathsf{ftargets}[1..n]$ (finalize targets, initially $\bot$).

\paragraph{processVote.} The finalize condition drops the $J \preceq v.\target$ guard and records the target:
\begin{center}
\begin{boxedminipage}{0.9\textwidth}
\begin{algorithmic}[1]
  \If{$v.\height = h_j$ and $v.\kind = \finalize$ and $\jtargets[h_j][i] \neq \bot$ and $v.\target \preceq \jtargets[h_j][i]$ and $v.\target \preceq B$}
    \State $P \gets P \cup \{i\}$
    \State $\mathsf{ftargets}[i] \gets v.\target$
  \EndIf
\end{algorithmic}
\end{boxedminipage}
\end{center}

\paragraph{processHeight.} The finality block ``$F \gets J$'' is replaced by a top-down walk over $\mathsf{ftargets}$:
\begin{center}
\begin{boxedminipage}{0.9\textwidth}
\begin{algorithmic}[1]
  \If{$|P| \geq 2n/3$} \Comment{Finality: walk finalize targets from $B$ toward genesis}
    \State $\cnt \gets 0$
    \For{$B'$ from $B$ down to $B_{\mathrm{genesis}}$}
      \State $\cnt \gets \cnt + |\{i : \mathsf{ftargets}[i] = B'\}|$
      \If{$\cnt \geq 2n/3$}
        \If{$B' \succ F$} $F \gets B'$ \EndIf \Comment{Only advance $F$}
        \State \textbf{break}
      \EndIf
    \EndFor
  \EndIf
\end{algorithmic}
\end{boxedminipage}
\end{center}
When $J$ advances (justification fires), $\mathsf{ftargets}$ is also deleted.

\section{Count notarizations top-down}

\label{alg:state-machine-topdown}
The following changes are applied to \Cref{alg:state-machine}.
$T^\mathsf{min}$ is removed from the state.

\paragraph{processVote.} $\ntargets[i]$ records the \emph{greatest} (longest) target.
The $T^\mathsf{min}$ update is removed.
\begin{center}
\begin{boxedminipage}{0.9\textwidth}
\begin{algorithmic}[1]
  \State \Comment{Both notarize and finalize votes}
  \If{$v.\target \succeq \ntargets[i]$}
    \State $\ntargets[i] \gets v.\target$
  \EndIf
\end{algorithmic}
\end{boxedminipage}
\end{center}

\paragraph{processHeight.} The notarization walk is replaced by a top-down walk from $B$ down to $T^\mathsf{max}$ (symmetric to the justification walk), counting $\ntargets[i]$:
\begin{center}
\begin{boxedminipage}{0.9\textwidth}
\begin{algorithmic}[1]
  \State $\cnt \gets 0$ \Comment{Notarization: walk from $B$ down to $T^\mathsf{max}$}
  \For{$B'$ from $B$ down to $T^\mathsf{max}$}
    \State $\cnt \gets \cnt + |\{i : \ntargets[i] = B'\}|$
    \If{$\cnt \geq 2n/3$} \Comment{Height advances}
      \If{$B' \succeq T^\mathsf{max}$}
        \State $T^\mathsf{max}\gets B'$
      \Else
        \State \textbf{delete} $\jtargets[h]$
      \EndIf
      \State $h \gets h + 1$
      \State \textbf{delete} $\ntargets$
      \State \Return $s$
    \EndIf
  \EndFor
\end{algorithmic}
\end{boxedminipage}
\end{center}

\begin{lemma}[Notarization respects finalization]
\Cref{alg:state-machine-topdown}.
Unless $\geq n/3$ validators are slashable: if $C$ is finalized at height~$h$ and block~$B$ is notarized at height~$h$ on some chain $X$, then $C \preceq X$.
\end{lemma}

\begin{proof}
Finalization of~$C$ at height~$h$ requires a quorum~$Q_F$ ($|Q_F| \geq 2n/3$) where each member~$i$ signed a finalize vote with $\height = h$ and target $D_i \succeq C$.
By slashing condition~2, no member of~$Q_F$ can have signed a notarize vote at height~$h$ without being slashable.
By slashing condition~1, each member has a unique finalize target~$D_i$.

Notarization of~$B$ at height~$h$ on chain~$X$ means the top-down notarize walk reached $2n/3$ at~$B$.
Let $Q_N = \{i : \ntargets[i] \succeq B\}$ on chain~$X$, with $|Q_N| \geq 2n/3$.
By quorum intersection, $|Q_F \cap Q_N| \geq n/3$.

Consider a non-slashable validator~$i$ in $Q_F \cap Q_N$.
Since~$i$ cannot have sent a notarize vote (condition~2), $\ntargets[i]$ on chain~$X$ was set by~$i$'s finalize vote.
For this vote to be processed on chain~$X$: $D_i \preceq X$.
From~$Q_F$: $D_i \succeq C$.
Hence $C \preceq D_i \preceq X$.
\end{proof}

\end{document}
